% GENERATED FILE. Edit canonical lessons, then run npm run generate:books.
% canonical-source-hash: fnv1a64:25499e15b349e5c3
% canonical-lessons: UR-C507-chakla, UR-C507-gundhna, UR-C507-belna, UR-C507-thapna, UR-C507-anch

\chapter{Making Rotis}
\label{ch:ur-making-rotis}
\hlchaptermodality{507}

\begin{chapteropening}
\textbf{By the end of this chapter:} \emph{I can say how bread is made at home.}

Complete the last lesson of chapter 507: I can say how bread is made at home.
\end{chapteropening}

\section[\texorpdfstring{\ur{چکلا} (chaklā)}{چکلا (chaklā)}]{\ur{چکلا} (chaklā) — a board for rolling dough}

\label{lesson:UR-C507-chakla}



\begin{quote}
Before the new one: say the Urdu for to be washed, then the Urdu for to rinse out.
\end{quote}

\subsection*{You'll want to know: \ur{چکلا}}
\textbf{\ur{چکلا}} — \emph{chaklā} — \textquotedblleft{}a board for rolling dough\textquotedblright{}.

\begin{quote}
\ur{چکلا} \ur{باورچی} \ur{خانے} \ur{میں} \ur{میز} \ur{پر} \ur{ہے۔}
\end{quote}

\begin{quote}
\emph{chaklā bāvarchī khāne meṅ mez par hai.}
\end{quote}

\begin{quote}
The rolling board is on the table in the kitchen.
\end{quote}

Masculine.

\begin{cousinweb}
\textbf{\ur{چکلا}} (\emph{chaklā}) is inherited: CDIAL lists the Hindi cousin \emph{caklā}, \textquotedblleft{}round and flat\textquotedblright{}, under \emph{*cakrala}, a form it reconstructs (CDIAL 4548).
\end{cousinweb}

\begin{grammarlens}[title={What you've built}]
One more word for describing what you lived and felt.
\end{grammarlens}

\subsection*{Your turn}
\begin{itemize}
\raggedright
  \item \emph{Say it:} \emph{chaklā}
  \item \emph{Say it:} \emph{chaklā}, once more
  \item \emph{Say it:} \emph{chaklā bāvarchī khāne meṅ mez par hai.} — The rolling board is on the table in the kitchen.
  \item \emph{Recall:} say the Urdu for to be washed, then the Urdu for to rinse out, then say \emph{chaklā} again
\end{itemize}

\begin{tcolorbox}[breakable,colback=teal!4,colframe=teal!35!black,title={Before you move on}]
What does \ur{چکلا} mean? (\textquotedblleft{}A board for rolling dough\textquotedblright{}.) Say it once more.
\end{tcolorbox}
\section[\texorpdfstring{\ur{گوندھنا} (gūṅdhnā)}{گوندھنا (gūṅdhnā)}]{\ur{گوندھنا} (gūṅdhnā) — to knead}

\label{lesson:UR-C507-gundhna}



\begin{quote}
Before the new one: say the Urdu for to rinse out, then the Urdu for a board for rolling dough.
\end{quote}

\subsection*{You'll want to know: \ur{گوندھنا}}
\textbf{\ur{گوندھنا}} — \emph{gūṅdhnā} — \textquotedblleft{}to knead\textquotedblright{}.

\begin{quote}
\ur{امی} \ur{صبح} \ur{آٹا} \ur{گوندھتی} \ur{ہیں۔}
\end{quote}

\begin{quote}
\emph{ammī subah āṭā gūṅdhtī haiṅ.}
\end{quote}

\begin{quote}
Mum kneads the dough in the morning.
\end{quote}

Past, with \ur{نے} (he or she …): \ur{اس} \ur{نے} \ur{آٹا} \ur{گوندھا}, \emph{us ne āṭā gūṅdhā}.

\begin{cousinweb}
\textbf{\ur{گوندھنا}} (\emph{gūṅdhnā}) is inherited: CDIAL lists the Hindi cousin \emph{gū̃dhnā}, \textquotedblleft{}to plait, to mix, to knead\textquotedblright{}, under Sanskrit \emph{guphati}, \textquotedblleft{}strings together\textquotedblright{} (CDIAL 4205).
\end{cousinweb}

\begin{grammarlens}[title={What you've built}]
One more word for describing what you lived and felt.
\end{grammarlens}

\subsection*{Your turn}
\begin{itemize}
\raggedright
  \item \emph{Say it:} \emph{gūṅdhnā}
  \item \emph{Say it:} \emph{gūṅdhnā}, once more
  \item \emph{Say it:} \emph{ammī subah āṭā gūṅdhtī haiṅ.} — Mum kneads the dough in the morning.
  \item \emph{Recall:} say the Urdu for to rinse out, then the Urdu for a board for rolling dough, then say \emph{gūṅdhnā} again
\end{itemize}

\begin{tcolorbox}[breakable,colback=teal!4,colframe=teal!35!black,title={Before you move on}]
What does \ur{گوندھنا} mean? (\textquotedblleft{}To knead\textquotedblright{}.) Say it once more.
\end{tcolorbox}
\section[\texorpdfstring{\ur{بیلنا} (belnā)}{بیلنا (belnā)}]{\ur{بیلنا} (belnā) — to roll out (dough)}

\label{lesson:UR-C507-belna}



\begin{quote}
Before the new one: say the Urdu for a board for rolling dough, then the Urdu for to knead.
\end{quote}

\subsection*{You'll want to know: \ur{بیلنا}}
\textbf{\ur{بیلنا}} — \emph{belnā} — \textquotedblleft{}to roll out (dough)\textquotedblright{}.

\begin{quote}
\ur{سارا} \ur{چکلے} \ur{پر} \ur{روٹی} \ur{بیل} \ur{رہی} \ur{تھی۔}
\end{quote}

\begin{quote}
\emph{Sārā chakle par roṭī bel rahī thī.}
\end{quote}

\begin{quote}
Sara was rolling out a roti on the board.
\end{quote}

Past, with \ur{نے} (he or she …): \ur{اس} \ur{نے} \ur{روٹی} \ur{بیلی}, \emph{us ne roṭī belī}.

\begin{cousinweb}
\textbf{\ur{بیلنا}} (\emph{belnā}) is inherited: CDIAL lists the Hindi cousin \emph{belnā}, \textquotedblleft{}to roll out\textquotedblright{}, under Sanskrit \emph{vēllayati} (CDIAL 12121).
\end{cousinweb}

\begin{grammarlens}[title={What you've built}]
One more word for describing what you lived and felt.
\end{grammarlens}

\subsection*{Your turn}
\begin{itemize}
\raggedright
  \item \emph{Say it:} \emph{belnā}
  \item \emph{Say it:} \emph{belnā}, once more
  \item \emph{Say it:} \emph{Sārā chakle par roṭī bel rahī thī.} — Sara was rolling out a roti on the board.
  \item \emph{Recall:} say the Urdu for a board for rolling dough, then the Urdu for to knead, then say \emph{belnā} again
\end{itemize}

\begin{tcolorbox}[breakable,colback=teal!4,colframe=teal!35!black,title={Before you move on}]
What does \ur{بیلنا} mean? (\textquotedblleft{}To roll out (dough)\textquotedblright{}.) Say it once more.
\end{tcolorbox}
\section[\texorpdfstring{\ur{تھاپنا} (thāpnā)}{تھاپنا (thāpnā)}]{\ur{تھاپنا} (thāpnā) — to pat (into shape)}

\label{lesson:UR-C507-thapna}



\begin{quote}
Before the new one: say the Urdu for to knead, then the Urdu for to roll out (dough).
\end{quote}

\subsection*{You'll want to know: \ur{تھاپنا}}
\textbf{\ur{تھاپنا}} — \emph{thāpnā} — \textquotedblleft{}to pat (into shape)\textquotedblright{}.

\begin{quote}
\ur{دادی} \ur{ہاتھوں} \ur{سے} \ur{مکئی} \ur{کی} \ur{روٹی} \ur{تھاپتی} \ur{ہیں۔}
\end{quote}

\begin{quote}
\emph{dādī hāthoṅ se makaī kī roṭī thāptī haiṅ.}
\end{quote}

\begin{quote}
Grandmother pats maize rotis into shape by hand.
\end{quote}

Past, with \ur{نے} (he or she …): \ur{اس} \ur{نے} \ur{روٹی} \ur{تھاپی}, \emph{us ne roṭī thāpī}.

\begin{cousinweb}
\textbf{\ur{تھاپنا}} (\emph{thāpnā}) is inherited: CDIAL lists the Hindi cousin \emph{thāpnā}, \textquotedblleft{}to pat\textquotedblright{}, under \emph{*thapp}, a form it reconstructs (CDIAL 6091).
\end{cousinweb}

\begin{grammarlens}[title={What you've built}]
One more word for describing what you lived and felt.
\end{grammarlens}

\subsection*{Your turn}
\begin{itemize}
\raggedright
  \item \emph{Say it:} \emph{thāpnā}
  \item \emph{Say it:} \emph{thāpnā}, once more
  \item \emph{Say it:} \emph{dādī hāthoṅ se makaī kī roṭī thāptī haiṅ.} — Grandmother pats maize rotis into shape by hand.
  \item \emph{Recall:} say the Urdu for to knead, then the Urdu for to roll out (dough), then say \emph{thāpnā} again
\end{itemize}

\begin{tcolorbox}[breakable,colback=teal!4,colframe=teal!35!black,title={Before you move on}]
What does \ur{تھاپنا} mean? (\textquotedblleft{}To pat (into shape)\textquotedblright{}.) Say it once more.
\end{tcolorbox}
\section[\texorpdfstring{\ur{آنچ} (āṅch)}{آنچ (āṅch)}]{\ur{آنچ} (āṅch) — the heat of a flame}

\label{lesson:UR-C507-anch}



\begin{quote}
Before the new one: say the Urdu for to roll out (dough), then the Urdu for to pat (into shape).
\end{quote}

\subsection*{You'll want to know: \ur{آنچ}}
\textbf{\ur{آنچ}} — \emph{āṅch} — \textquotedblleft{}the heat of a flame\textquotedblright{}.

\begin{quote}
\ur{دال} \ur{کو} \ur{دھیمی} \ur{آنچ} \ur{پر} \ur{پکائیے۔}
\end{quote}

\begin{quote}
\emph{dāl ko dhīmī āṅch par pakāiye.}
\end{quote}

\begin{quote}
Cook the lentils on a low flame, please.
\end{quote}

Feminine.

\begin{cousinweb}
\textbf{\ur{آنچ}} (\emph{āṅch}) is inherited from Sanskrit \emph{arcis}, \textquotedblleft{}a flame, a ray\textquotedblright{} (CDIAL 635).
\end{cousinweb}

\begin{grammarlens}[title={What you've built}]
One more word for describing what you lived and felt.
\end{grammarlens}

\subsection*{Your turn}
\begin{itemize}
\raggedright
  \item \emph{Say it:} \emph{āṅch}
  \item \emph{Say it:} \emph{āṅch}, once more
  \item \emph{Say it:} \emph{dāl ko dhīmī āṅch par pakāiye.} — Cook the lentils on a low flame, please.
  \item \emph{Recall:} say the Urdu for to roll out (dough), then the Urdu for to pat (into shape), then say \emph{āṅch} again
\end{itemize}

\begin{tcolorbox}[breakable,colback=teal!4,colframe=teal!35!black,title={Before you move on}]
What does \ur{آنچ} mean? (\textquotedblleft{}The heat of a flame\textquotedblright{}.) Say it once more.
\end{tcolorbox}
